% Please use LuaLaTeX or XeLaTeX for compilation
\documentclass[11pt,aspectratio=169]{beamer}

% Theme and Colors
\usetheme{eth}
\colorlet{titlefgcolor}{ETHGray}
\colorlet{accentcolor}{ETHRed}
\colorlet{titlebgcolor}{ETHRed} % Title slide background

% Standard Packages
\usepackage{amsmath, amssymb} % Math symbols
\usepackage{bm} % Bold math symbols
\usepackage{graphicx} % Include graphics (if needed)
\usepackage{booktabs} % Nicer tables (if needed)
\usepackage{microtype} % Better typography
\usepackage{physics} % Physics notation

% Ben packages
\usepackage{soul}

% Citation Management (using biblatex)
\usepackage[style=numeric, backend=biber]{biblatex}
\addbibresource{references.bib} % You need a references.bib file

% Footnotes (optional, if needed)
\usepackage[multiple]{footmisc}
\usepackage{footnoterange}
\usepackage{fnpct}
\AdaptNote\footcite{oo+m}[footnote]{%
\setfnpct{dont-mess-around}%
\IfNoValueTF{#1}
{#NOTE{#3}}
{\IfNoValueTF{#2}{#NOTE[#1]{#3}}{#NOTE[#1][#2]{#3}}}%
}

% Algorithms (optional, if needed)
\usepackage{algorithm}
\usepackage[noend]{algpseudocode} % Use noend option as in template

% Tikz/PGFPlots (optional, if needed for diagrams)
\usepackage{tikz}
\usetikzlibrary{shapes, arrows, positioning, calc}
\usepackage{pgfplots}
\pgfplotsset{compat=1.17} % Use a recent compatibility level

% Custom command for highlighting math
\newcommand\mathalert[1]{{\color{accentcolor}#1}}

% --- Presentation Metadata ---
\title{The Berezinskii-Kosterlitz-Thouless Transition in the 2D XY Model}
\subtitle{Monte Carlo Simulation and Analysis}
\author{Lara Turgut \and Francesco Conoscenti \and Ben Bullinger}
\institute{ETH Zürich\\ Department of Physics}
\date{Computational Statistical Physics (402-0812-00L), Spring 2025}

\begin{document}

% --- Title Slide ---
\def\titlefigure{} % No figure on title slide unless specified
\setlength{\titleboxwidth}{0.85\textwidth} % Adjust box width if needed
\titleframe % Creates the title slide using the theme's definition

% --- Table of Contents ---
\tocframe % Creates the table of contents slide

% --- Section 1: Introduction ---
\section{Introduction}

\begin{frame}{The 2D XY Model}
    \begin{itemize}
        \item XY model: spins constrained to a plane with $O(2)$ symmetry
        \begin{equation}
            H = -J \sum_{\langle i,j \rangle} \mathbf{S}_i \cdot \mathbf{S}_j = -J \sum_{\langle i,j \rangle} \cos(\theta_i - \theta_j)
        \end{equation}
        \item Each spin is a unit vector $\mathbf{S}_i = (\cos\theta_i, \sin\theta_i)$
        \pause
        \item Unlike conventional phase transitions with symmetry breaking:
        \begin{itemize}
            \item \textbf{Mermin-Wagner theorem:} No spontaneous symmetry breaking (true long-range order) at $T > 0$ in systems with continuous symmetry in $d \leq 2$
            \item Yet, the 2D XY model \textit{does} exhibit a phase transition!
        \end{itemize}
        \pause
        \item The Berezinskii-Kosterlitz-Thouless (BKT) transition:
        \begin{itemize}
            \item Topological phase transition with no conventional order parameter
            \item Transition from quasi-long-range order to disorder
            \item Driven by the unbinding of vortex-antivortex pairs
        \end{itemize}
    \end{itemize}
\end{frame}

\begin{frame}{Physical Significance}
    \begin{itemize}
        \item The BKT transition represents a unique phase transition mechanism
        \begin{itemize}
            \item \textbf{Infinite-order transition:} All derivatives of free energy are continuous
            \item \textbf{Topological in nature:} Driven by topological defects (vortices)
        \end{itemize}
        \pause
        \item Found in diverse physical systems:
        \begin{itemize}
            \item Thin films of superfluid helium-4
            \item 2D superconductors
            \item 2D melting (KTHNY theory)
            \item Josephson junction arrays
            \item Cold atom systems
            \item Quantum spin systems
        \end{itemize}
        \pause
        \item Universal mechanism: A beautiful example of universality in critical phenomena
    \end{itemize}
\end{frame}

% --- Section 2: Theoretical Background ---
\section{Theoretical Background}

\begin{frame}{Algebraic Order in the 2D XY Model}
    \begin{itemize}
        \item Low-temperature expansion of the Hamiltonian:
        \begin{equation}
            H \approx E_0 + \frac{J}{2} \int d^2r (\nabla \theta(\mathbf{r}))^2
        \end{equation}
        \pause
        \item Correlation function analysis shows:
        \begin{equation}
            g(r) = \langle e^{i(\theta(\mathbf{r}) - \theta(0))} \rangle \approx \left( \frac{\pi r}{a} \right)^{-\frac{k_B T}{2\pi J}}
        \end{equation}
        \item Key property: \mathalert{algebraic decay} of correlations at low temperature
        \begin{itemize}
            \item \textbf{Low-T phase:} $g(r) \sim r^{-\eta(T)}$ with temperature-dependent exponent
            \item \textbf{High-T phase:} $g(r) \sim e^{-r/\xi(T)}$ (exponential decay)
        \end{itemize}
        \pause
        \item This algebraic decay = \textbf{quasi-long-range order (QLRO)}
        \begin{itemize}
            \item Infinite correlation length $\xi = \infty$ in low-T phase
            \item System is critical at all temperatures below $T_{KT}$
        \end{itemize}
    \end{itemize}
\end{frame}

\begin{frame}{The Role of Vortices}
    \begin{itemize}
        \item Vortices: topological defects with winding of $2\pi n$ in phase field
        \begin{equation}
            \oint \nabla \theta \cdot d\mathbf{l} = 2\pi n
        \end{equation}
        \item Energy of a single vortex (with $n=1$):
        \begin{equation}
            \Delta E = \pi J \ln \frac{L}{a}
        \end{equation}
        where $L$ is system size and $a$ is lattice spacing
        \pause
        \item Entropy of a single vortex:
        \begin{equation}
            \Delta S = k_B \ln \frac{L^2}{a^2} = 2k_B \ln \frac{L}{a}
        \end{equation}
        \pause
        \item Free energy:
        \begin{equation}
            \Delta F = \Delta E - T\Delta S = (\pi J - 2k_B T)\ln \frac{L}{a}
        \end{equation}
        \pause
        \item \mathalert{Crucial insight}: Vortices are bound in pairs below critical temperature, and unbound above it
        \begin{itemize}
            \item For $T < \frac{\pi J}{2k_B}$: $\Delta F > 0$ (bound vortex pairs)
            \item For $T > \frac{\pi J}{2k_B}$: $\Delta F < 0$ (free vortices proliferate)
        \end{itemize}
    \end{itemize}
\end{frame}

\begin{frame}{Renormalization Group Analysis}
    \begin{itemize}
        \item Mapping to 2D Coulomb gas: vortices as charges with $\ln(r)$ interaction
        \item RG flow equations for coupling $K = \frac{J}{k_B T}$ and fugacity $y$:
        \begin{align}
            \frac{d K^{-1}(l)}{dl} &= 4 \pi^3 y(l)^2 + O(y(l)^4) \\
            \frac{d y(l)}{dl} &= (2 - \pi K(l)) y(l) + O(y(l)^3)
        \end{align}
        \pause
        \item Flow diagram shows two phases:
        \begin{itemize}
            \item \textbf{Low-T} ($K > \frac{2}{\pi}$): $y \to 0$ (no free vortices, QLRO)
            \item \textbf{High-T} ($K < \frac{2}{\pi}$): $y$ increases (free vortices, disorder)
        \end{itemize}
        \pause
        \item Universal properties:
        \begin{itemize}
            \item Critical point: $K_c = \frac{2}{\pi}$ or $T_{KT} = \frac{\pi J}{2k_B}$
            \item \textbf{Nelson-Kosterlitz jump}: Spin stiffness $\rho_s$ jumps at transition:
            \begin{equation}
                \rho_s(T_{KT}^-) = \frac{2T_{KT}}{\pi}
            \end{equation}
            \item Correlation length above transition: $\xi(T) \sim e^{b/\sqrt{T/T_{KT}-1}}$
        \end{itemize}
    \end{itemize}
\end{frame}

% --- Section 3: Monte Carlo Methods ---
\section{Monte Carlo Methods}

\begin{frame}{Simulating the XY Model}
    \begin{itemize}
        \item Challenges in simulating BKT transition:
        \begin{itemize}
            \item Critical slowing down near phase transition
            \item Need for large system sizes (finite-size effects)
            \item Long correlation times
        \end{itemize}
        \pause
        \item Standard Metropolis algorithm:
        \begin{itemize}
            \item Energy change for spin flip: $\Delta E = -J \sum_{\langle i,j \rangle} [\cos(\theta'_i - \theta_j) - \cos(\theta_i - \theta_j)]$
            \item Accept new state with probability $P = \min(1, e^{-\beta \Delta E})$
            \item Very inefficient near critical point due to large correlated regions
        \end{itemize}
        \pause
        \item More efficient approaches:
        \begin{itemize}
            \item \textbf{Wolff cluster algorithm}: Updates entire clusters of spins
            \item \textbf{Parallel tempering}: Exchanges configurations between different temperatures
            \item Combination of both methods for optimal efficiency
        \end{itemize}
    \end{itemize}
\end{frame}

\begin{frame}{Wolff Cluster Algorithm}
    \begin{itemize}
        \item Efficiently updates entire clusters of correlated spins
        \begin{enumerate}
            \item Choose a random reflection axis (unit vector $\mathbf{r}$)
            \item Pick a random initial site for the cluster
            \item Add neighboring spins with probability $P_{\text{add}} = 1 - \exp[\min(0, 2\beta J (\mathbf{S}_i \cdot \mathbf{r})(\mathbf{S}_j \cdot \mathbf{r}))]$
            \item Reflect all spins in the cluster w.r.t. the chosen axis
        \end{enumerate}
        \pause
        \item Advantages:
        \begin{itemize}
            \item Dramatically reduces critical slowing down
            \item Dynamical critical exponent $z \approx 0$ (vs. $z \approx 2$ for Metropolis)
            \item Efficiently samples low-energy configurations
        \end{itemize}
        \pause
        \item Our implementation:
        \begin{itemize}
            \item Combines Wolff updates with parallel tempering
            \item Optimally adapts the number of updates based on average cluster size
            \item Uses numba JIT compilation for performance
        \end{itemize}
    \end{itemize}
\end{frame}

\begin{frame}{Parallel Tempering}
    \begin{itemize}
        \item Simulates multiple replicas at different temperatures $T_1 < T_2 < ... < T_M$
        \item Periodically proposes exchanges between adjacent temperatures
        \begin{itemize}
            \item Accept exchange with probability $P = \min(1, \exp((\beta_i - \beta_j)(E_j - E_i)))$
        \end{itemize}
        \pause
        \item Benefits:
        \begin{itemize}
            \item Helps systems escape from local minima
            \item Improves sampling at low temperatures
            \item Configurations can diffuse across temperature space
        \end{itemize}
        \pause
        \item Implementation details:
        \begin{itemize}
            \item Alternating even/odd exchanges for better mixing
            \item Temperature spacing optimized for exchange acceptance rates
            \item Combined with Wolff algorithm for optimal efficiency
        \end{itemize}
    \end{itemize}
\end{frame}

% --- Section 4: Implementation ---
\section{Implementation Details}

\begin{frame}{Our Implementation Approach}
    \begin{itemize}
        \item Three-phase simulation structure:
        \begin{enumerate}
            \item \textbf{Pre-thermalization}: Determine optimal number of Wolff updates
            \item \textbf{Thermalization}: Reach equilibrium at each temperature
            \item \textbf{Measurement}: Collect data for analysis
        \end{enumerate}
        \pause
        \item Optimization techniques:
        \begin{itemize}
            \item Numba JIT compilation for compute-intensive functions
            \item Vectorized operations with NumPy for efficiency
            \item Pre-computed neighbor lists to avoid redundant calculations
            \item Dynamic adjustment of updates based on cluster sizes
        \end{itemize}
        \pause
        \item Simulation flow:
        \begin{itemize}
            \item Generate temperature range (linear or geometric spacing)
            \item Run simulations for multiple system sizes $L$
            \item For each size and temperature: thermalize, measure, analyze
            \item Perform finite-size scaling analysis to extract $T_{KT}$ 
        \end{itemize}
    \end{itemize}
\end{frame}

\begin{frame}{Key Observables}
    \begin{itemize}
        \item Physical quantities measured in our simulation:
        \begin{itemize}
            \item \textbf{Energy per spin}: $e = \frac{\langle H \rangle}{N} = -\frac{J}{N} \sum_{\langle i,j \rangle} \langle\cos(\theta_i - \theta_j)\rangle$
            \item \textbf{Specific heat}: $c_v = \frac{\langle E^2 \rangle - \langle E \rangle^2}{NT^2}$
            \item \textbf{Magnetization}: $\mathbf{m} = \frac{1}{N}\left\langle \sum_i (\cos\theta_i, \sin\theta_i) \right\rangle$
            \item \textbf{Susceptibility}: $\chi = \frac{\langle M^2 \rangle - \langle M \rangle^2}{NT}$
            \pause
            \item \textbf{Spin stiffness (Helicity modulus)}:
            \begin{equation}
                \rho_s = \langle H_x \rangle - \frac{\beta}{N} \langle I_x^2 \rangle
            \end{equation}
            where $H_x = \sum_{\langle i,j \rangle_x} \cos(\theta_i - \theta_j)$ and $I_x = \sum_{\langle i,j \rangle_x} \sin(\theta_i - \theta_j)$
            \pause
            \item \textbf{Vortex density}: $\omega_v = \frac{1}{N} \sum_{\text{plaquettes}} \left| \frac{1}{2\pi} \sum_{\text{edges}} \Delta\theta_{ij} \right|$
        \end{itemize}
        \item Spin stiffness is crucial for locating $T_{KT}$ through the universal jump
    \end{itemize}
\end{frame}

\begin{frame}{Error Analysis}
    \begin{itemize}
        \item Statistical errors from Monte Carlo sampling:
        \begin{itemize}
            \item Sequential MC data is \textit{correlated}
            \item Standard error formulas for independent data are inadequate
        \end{itemize}
        \pause
        \item Jackknife resampling for reliable error estimation:
        \begin{enumerate}
            \item Divide data into $B$ blocks larger than correlation time
            \item For each block $b$, compute estimate $\bar{O}_b$ by excluding block $b$
            \item Jackknife error: $\sigma^2_{\text{JK}} = \frac{B-1}{B}\sum_b (\bar{O}_b - \bar{O})^2$
        \end{enumerate}
        \pause
        \item Benefits:
        \begin{itemize}
            \item Correctly accounts for autocorrelations in MC data
            \item Works well for derived quantities (non-linear functions)
            \item More reliable than naive variance calculations
            \item Essential for accurate error bars on quantities near transition
        \end{itemize}
    \end{itemize}
\end{frame}

% --- Section 5: Results and Analysis ---
\section{Results and Analysis}

\begin{frame}{Thermodynamic Observables}
    \begin{itemize}
        \item Specific heat $c_v$:
        \begin{itemize}
            \item Broad peak near $T \approx 1.0 J$
            \item No divergence at $T_{KT}$ (characteristic of BKT transition)
            \item Slight sharpening with increasing system size
        \end{itemize}
        \pause
        \item Magnetization:
        \begin{itemize}
            \item Decreases with system size at all temperatures
            \item Follows $m \sim L^{-\eta/2}$ in low-temperature phase (QLRO)
            \item Approaches zero exponentially fast with $L$ above $T_{KT}$
        \end{itemize}
        \pause
        \item Susceptibility:
        \begin{itemize}
            \item Peaks near transition temperature
            \item Peak height increases with system size
            \item Shows logarithmic divergence in QLRO phase
        \end{itemize}
    \end{itemize}
\end{frame}

\begin{frame}{Identifying the BKT Transition}
    \begin{itemize}
        \item Vortex density:
        \begin{itemize}
            \item Very low at $T \ll T_{KT}$ (few vortex pairs)
            \item Increases sharply around $T \approx 0.9 J$ (vortex unbinding)
            \item Provides visual indication of transition
        \end{itemize}
        \pause
        \item Spin stiffness (helicity modulus):
        \begin{itemize}
            \item Key indicator of the BKT transition
            \item Universal jump at $T_{KT}$: $\rho_s(T_{KT}^-) = \frac{2T_{KT}}{\pi}$
            \item Intersection of $\rho_s(T)$ curves with the line $\rho_s = \frac{2T}{\pi}$ gives $T_{KT}$
            \item Our data indicates $T_{KT} \approx 0.89 J$
        \end{itemize}
        \pause
        \item Finite-size effects:
        \begin{itemize}
            \item Transition is smoothed out in finite systems
            \item Slopes become sharper with increasing $L$
            \item True universal jump only emerges in thermodynamic limit $L \to \infty$
        \end{itemize}
    \end{itemize}
\end{frame}

\begin{frame}{Finite-Size Scaling}
    \begin{itemize}
        \item Theory predicts finite-size scaling of $T_{KT}(L)$:
        \begin{equation}
            T_{KT}(L) = T_{KT}(\infty) + \frac{a}{(\ln L)^2}
        \end{equation}
        \pause
        \item Extracting $T_{KT}(L)$ for each system size:
        \begin{itemize}
            \item For each $L$, find temperature where $\rho_s(T) = 2T/\pi$
            \item Extract $T_{KT}(L)$ and its uncertainty
        \end{itemize}
        \pause
        \item Finite-size scaling analysis:
        \begin{itemize}
            \item Plot $T_{KT}(L)$ vs. $1/(\ln L)^2$
            \item Linear fit gives $T_{KT}(\infty)$ (thermodynamic limit)
            \item Our analysis yields $T_{KT}(\infty) \approx 0.893(4) J$
            \item Consistent with best theoretical and numerical estimates in literature
        \end{itemize}
    \end{itemize}
\end{frame}

% --- Section 6: Conclusion ---
\section{Conclusion}

\begin{frame}{Summary and Significance}
    \begin{itemize}
        \item Successfully simulated the BKT transition in the 2D XY model:
        \begin{itemize}
            \item Implemented efficient Wolff algorithm with parallel tempering
            \item Measured key observables across multiple system sizes
            \item Applied finite-size scaling to extract critical temperature
        \end{itemize}
        \pause
        \item Key findings:
        \begin{itemize}
            \item Confirmed the characteristic signatures of the BKT transition
            \item Determined $T_{KT} \approx 0.893(4) J$ in good agreement with theory
            \item Observed universal jump in spin stiffness
            \item Demonstrated vortex unbinding mechanism
        \end{itemize}
        \pause
        \item Broader implications:
        \begin{itemize}
            \item BKT transition represents a unique universality class
            \item Demonstrates the power of topological defects in phase transitions
            \item Provides insight into universal mechanisms across diverse physical systems
            \item Highlights importance of advanced sampling techniques for critical phenomena
        \end{itemize}
    \end{itemize}
\end{frame}

\begin{frame}{Future Directions}
    \begin{itemize}
        \item Methodological improvements:
        \begin{itemize}
            \item Further optimization for larger system sizes
            \item GPU implementation for parallel computation
            \item Advanced analysis techniques for critical properties
        \end{itemize}
        \pause
        \item Extensions to related models:
        \begin{itemize}
            \item XY model with disorder or frustration
            \item Quantum XY models and quantum phase transitions
            \item Higher-dimensional models with topological transitions
        \end{itemize}
        \pause
        \item Applications:
        \begin{itemize}
            \item Superconducting thin films and arrays
            \item Quantum computing platforms based on 2D systems
            \item Cold atom experiments with vortex dynamics
            \item Theoretical connections to fundamental physics
        \end{itemize}
    \end{itemize}
\end{frame}

% --- Bibliography ---
\begin{frame}[allowframebreaks]{References}
    % Create a references.bib file with appropriate entries
    \begin{thebibliography}{9}
        \bibitem{Berezinskii} V. L. Berezinskii, \textit{Sov. Phys. JETP} \textbf{32}, 493 (1971)
        \bibitem{KT} J. M. Kosterlitz and D. J. Thouless, \textit{J. Phys. C} \textbf{6}, 1181 (1973)
        \bibitem{Wolff} U. Wolff, \textit{Phys. Rev. Lett.} \textbf{62}, 361 (1989)
        \bibitem{JKKN} J. V. José et al., \textit{Phys. Rev. B} \textbf{16}, 1217 (1977)
        \bibitem{Nelson} D. R. Nelson and J. M. Kosterlitz, \textit{Phys. Rev. Lett.} \textbf{39}, 1201 (1977)
        \bibitem{BishopReppy} D. J. Bishop and J. D. Reppy, \textit{Phys. Rev. Lett.} \textbf{40}, 1727 (1978)
        \bibitem{KTHNY} B. I. Halperin and D. R. Nelson, \textit{Phys. Rev. Lett.} \textbf{41}, 121 (1978)
    \end{thebibliography}
\end{frame}

% --- Closing Slide ---
\begin{closingframe}
    \begin{center}
        \begin{minipage}{1\textwidth}
            \textbf{Thank you for your attention}
            
            \vspace{0.5em}
            
            \textit{Questions?}
        \end{minipage}
    \end{center}

    \vfill

    \begin{columns}
		\begin{column}{\textwidth}
            ETH Zurich\\
            Department of Physics
		\end{column}
	\end{columns}
\end{closingframe}

\end{document}

% Create a references.bib file with appropriate entries:
/*
@article{Berezinskii,
  title={Destruction of long-range order in one-dimensional and two-dimensional systems having a continuous symmetry group},
  author={Berezinskii, V L},
  journal={Soviet Physics JETP},
  volume={32},
  number={3},
  pages={493--500},
  year={1971}
}

@article{KT,
  title={Ordering, metastability and phase transitions in two-dimensional systems},
  author={Kosterlitz, J M and Thouless, D J},
  journal={Journal of Physics C: Solid State Physics},
  volume={6},
  number={7},
  pages={1181},
  year={1973},
  publisher={IOP Publishing}
}

@article{Wolff,
  title={Collective Monte Carlo updating for spin systems},
  author={Wolff, Ulli},
  journal={Physical Review Letters},
  volume={62},
  number={4},
  pages={361},
  year={1989},
  publisher={APS}
}

@article{JKKN,
  title={Renormalization, vortices, and symmetry-breaking perturbations in the two-dimensional planar model},
  author={Jos{\'e}, Jorge V and Kadanoff, Leo P and Kirkpatrick, Scott and Nelson, David R},
  journal={Physical Review B},
  volume={16},
  number={3},
  pages={1217},
  year={1977},
  publisher={APS}
}

@article{Nelson,
  title={Universal jump in the superfluid density of two-dimensional superfluids},
  author={Nelson, David R and Kosterlitz, John M},
  journal={Physical Review Letters},
  volume={39},
  number={19},
  pages={1201},
  year={1977},
  publisher={APS}
}

@article{BishopReppy,
  title={Study of the Superfluid Transition in Two-Dimensional $^4$He Films},
  author={Bishop, David J and Reppy, John D},
  journal={Physical Review Letters},
  volume={40},
  number={26},
  pages={1727},
  year={1978},
  publisher={APS}
}

@article{KTHNY,
  title={Phase transitions in two-dimensional systems},
  author={Halperin, B I and Nelson, David R},
  journal={Physical Review Letters},
  volume={41},
  number={2},
  pages={121},
  year={1978},
  publisher={APS}
}
*/
